\section{Polarimetric products (COV \& GCOV)}
\label{sec:pol_products}

\subsection{Product Description}
\label{sec:description_pol}

\acs{sar} polarimetry entails the simultaneous acquisition of a set of \acs{slc} images using different combinations of transmitted and received polarisations, and each combination is denoted as \textit{polarimetric channel}. Depending on the sensor and acquisition mode, there may be only one polarimetric channel (single-pol mode), two channels (dual-pol mode), or four channels (quad-pol mode). Using as inputs the \acs{slc} images at the polarimetric channels available, the polarimetric product is formed as a covariance matrix, as described in D1~\cite{D1}. Therefore, this \acs{ard} product contains the entries of the polarimetric covariance matrix, either in radar coordinates (\acs{cov} product) or in cartographical coordinates (\acs{gcov} product). In order to preserve the polarimetric information, it is necessary to apply the same processing to generate each one of the matrix entries. All the possible forms of the covariance matrix are detailed in D1~\cite{D1}, and not repeated here to save space.

A convenient way to explain the generation and the resulting format of the \acs{cov} and \acs{gcov} products is by re-using what is explained in \autoref{s:nrb} for the \acs{nrb} product. All matrix entries are processed as if each one of them constitutes a \acs{nrb} product. For instance, they
usually undergo Radiometric Terrain Correction (\acs{rtc}), providing Gamma-Nought as defined in \cite{Small2011}, which mitigates the impact of differing observation geometries. 

In fact, the diagonal entries of the covariance matrix correspond to the backscattering coefficient of each individual polarimetric channel, so the diagonal entries in a \acs{gcov} product coincide with the \acs{nrb} products of all channels. For instance, taking the VV-VH dual-pol mode as an example, typical of Sentinel-1 acquisitions over land, the diagonal elements  correspond to the \acs{nrb} of the VV and VH channels. As for the off-diagonal entries, they are formed by applying the same processing of the \acs{nrb} generation but to the product VV$\cdot$VH$^*$, where $*$ denotes the complex conjugate, instead of to the intensity of the VV and VH channels. In summary, \acs{cov} and \acs{gcov} products for the dual-pol case ($C_2$ matrix) contain only three layers: the two real-valued intensities (backscattering coefficients) of the two channels, and one complex value corresponding to the correlation between the two channels.

It must be pointed out that the \acs{nrb} product may be regarded either as a \textit{subset} or as a \textit{particular case} of the \acs{gcov} product, as it is fully contained in it. Such a point of view is followed by NASA in the \acs{gcov} product defined for NISAR~\cite{NISAR_GCOV_ProdSpec}, in which the \acs{nrb} product is just a particular case included in the \acs{gcov} specification.

The \acs{gcov} product shares all the properties of \acs{nrb} products, hence  enabling direct visualisation in geographical information systems for interpretation and analysis, but enriches the information content of \acs{nrb} products with the extra channels. Polarimetry outperforms backscatter in most final applications, as demonstrated in literature~\cite{Hajnsek2021}. 

Regarding the \acs{cov} product, i.e. expressed in radar coordinates, it may be preferred by some users which need to apply some specific processing steps in radar geometry to avoid distortions in the information content, like speckle filtering with a multi-look or boxcar, before geocoding the resulting matrix. In such a case, the \acs{cov} product can be provided together with a look-up table correspondence for geocoding, to be applied later by the user to convert the product to cartographical coordinates. 

As in the case of \acs{nrb}, the \acs{cov} and \acs{gcov} products can be obtained through two possible alternative processing flows:
\begin{enumerate}
    \item Starting from \acs{slc} data with the same Data Take ID, subject to the co-registration procedure (\acs{cslc});
    \item Starting from \acs{slc} data with the same Data Take ID, without applying the co-registration procedure.
\end{enumerate}

Due to the specific nature of the polarimetric \acs{ard} products, some considerations and clarifications are detailed next. 
%In addition, as for other products, existing SW tools and computational complexity aspects are commented in \autoref{ss:algoritms_pol_SW} and \autoref{ss:algoritms_pol_comput}, respectively.

\paragraph*{Reference channels}

Some processing associated with the co-registration and geocoding steps may be carried out using the actual data from the images. For instance, co-registration based on enhanced spectral diversity (ESD) exploits areas imaged by overlapping bursts. Similarly, point targets can be used to refine azimuth times of a stack of images with respect to a reference image. Regarding geocoding, a synthetic or simulated image generated from a \acs{dem} can be cross-correlated with the \acs{sar} image to improve the geocoding accuracy.
In all these cases, the computations (e.g. cross-correlation)  must be performed using only ONE of the polarimetric channels, whereas the result (e.g. interpolation of the secondary images at the computed positions) is applied to ALL the polarimetric channels.

For these cases, the selected channel is usually a co-polar one, i.e. HH or VV, because it provides a higher backscatter level and hence a higher \acs{snr} than the cross-polar channels. Consequently, the reference channel must be indicated in the metadata.

\paragraph*{Polarimetric symmetrization}

Every image acquired in quad-polarisation mode consists of four \acs{slc} images measured at the following channels: HH, HV, VH, and VV. Due to reciprocity, and after an adequate calibration, the two cross-polar channels (HV and VH) should be equal and are only different due to the presence of noise. For this reason, for the definition of the covariance matrix, only three channels are considered, i.e. HH, VV and one cross-polar channel (HV or VH). 

A common practice, also considered in the definition of the \acs{gcov}-NISAR product~\cite{NISAR_GCOV_ProdSpec}, consists in averaging the two cross-polar channels and using the resulting average in the construction of the covariance matrix in the \acs{ard} product. This operation is called symmetrization in polarimetric literature.

This step is not commented by \acs{ceos}, but the \acs{gcov}-NISAR definition in~\cite{NISAR_GCOV_ProdSpec} indicates that the metadata should include:
\begin{itemize}
\item A flag to indicate whether the symmetrization has been applied or not.
\item A reference to the symmetrization algorithm employed, if applicable.
\end{itemize}

\paragraph*{Order of spatial filtering and geocoding}
\label{p:order_pol}

In order to minimise the bias in the \acs{gcov} products, speckle filtering and/or spatial averaging of the covariance matrix should be applied in radar coordinates (slant-range and cross-range). Afterwards, with a resolution degraded by this filtering, geocoding should be performed using either nearest-neighbour resampling or the approach proposed in~\cite{Shiroma2022}. Specifically, nearest-neighbour resampling ensures that the averaged covariance
matrix elements in slant range and in geocoded ground projection are exactly the same (obviously assuming similar posting in both coordinate projections). If the approach in~\cite{Shiroma2022} is used, the degraded resolution makes it equivalent to nearest-neighbour resampling except in the borders between pixels, but no significant artifact is expected.

In the definition of the \acs{gcov}-NISAR product~\cite{NISAR_GCOV_ProdSpec}, the estimation of the output covariance matrix and its projection onto geographical coordinates is carried out simultaneously by means of the area-based interpolation associated with the \acs{rtc} proposed by Shiroma et al. in~\cite{Shiroma2022}. Therefore, in that case there is no degradation of the polarimetric information, since the averaging is carried out using the entries of the covariance matrix of the original pixels in radar coordinates.

\paragraph*{Polarimetric speckle filtering}

Speckle filter is required for extracting useful quantitative information from the \acs{sar} data, especially over natural areas. It is usually carried out using a spatial filtering of the image, which entails a degradation of the spatial resolution. To mitigate this issue, spatially-adaptive speckle filters are employed, which try to simultaneously preserve details and average homogeneous distributed areas. For this purpose, at each pixel to be filtered, a group of pixels with similar scattering properties is selected in an adaptive fashion. In this context, spatially-adaptive speckle filtering can clearly benefit from the polarimetric information of the images, expressed in covariance matrix form, since it allows for identifying similar and different land covers and elements in the scene. 

Many speckle filtering algorithms have been proposed which incorporate the polarimetric information in the selection of the group of homogeneous pixels, like the polarimetric intensity-driven adaptive-neighbourhood (IDAN) filter~\cite{Vasile2006}, the non-local filter~\cite{Deledalle2015}, tensor decomposition-based filter~\cite{Luo2023}, etc.
In all cases, the obtained performance is better than the one obtained with filters which do not exploit this information or with filters with fixed kernels (e.g. boxcar).

Speckle filters should be applied in radar coordinates (i.e. before geocoding). In this sense, the simultaneous area-based averaging and geocoding considered in the \acs{gcov}-NISAR product~\cite{NISAR_GCOV_ProdSpec} disables the possibility to rigorously apply these adaptive filters to the (already geocoded) product provided as a \acs{gcov}. It is for this reason that the \acs{cov} product is required by some users.


\paragraph*{Considerations about the Equivalent Number of Looks}

The use of spatially-adaptive speckle filters and/or the area-based interpolation associated with the \acs{rtc} defined by Shiroma et al. in~\cite{Shiroma2022} yields a variable equivalent number of looks (ENL) over the image. As a result, the layer with the values of ENL has to be available in the per-pixel metadata of the \acs{ard} products.

The ENL value is later used to help interpret or exploit some derived products. For instance, the maximum likelihood change detector for polarimetric images (also valid for single-pol backscatter data) is based on the generalised likelihood ratio test computed by using covariance matrices~\cite{Conradsen2003, Nielsen2015}. In this test, many parameters (statistical significance, false alarm probability, thresholds, etc.) depend on the ENL of the two pixels being compared, so this information must be provided together with the \acs{ard} product (as a metadata layer).

Another aspect to highlight regarding ENL is the minimum required value to properly employ some techniques. The multi-dimensional nature of the polarimetric information, exemplified by the matrix form of the data (i.e. the covariance matrix) makes it necessary to have a large number of looks to measure some polarimetric features. For instance, the entropy, alpha and anisotropy products~\cite{Cloude1997, Lee_Pottier2009} derived from the eigenvector decomposition of the covariance matrix (after transformation to coherency matrix) are widely employed in scientific studies. The estimation of these products is known to be biased when the number of looks is low, requiring large values (50+) for sufficiently reducing that bias~\cite{LopezMartinez2005, Lee2008}.

As a consequence, the ENL resulting from the area-based interpolation associated with the \acs{rtc} defined by Shiroma et al. in~\cite{Shiroma2022}, which is used in the \acs{gcov}-NISAR product~\cite{NISAR_GCOV_ProdSpec}, is likely to be smaller than the one required for some applications. The same will happen when the speckle filter applied before the geocoding of the \acs{pol} product defined in \acs{ceos}~\cite{CEOS-ARD_PFS_SAR} employs a small value of looks (e.g. a 3x3 or 5x5 boxcar). In both cases, the usually adopted solution consists in adding an extra step of spatial averaging to the previously-filtered covariance matrices and before the computation of the derived products. This extra filtering step, if applied after geocoding, will produce some artifacts as it was commented in \autoref{p:order_pol}.

The severity of the impact of these artifacts will depend on the specific case study and on the final application, but it must be taken into account.

\paragraph*{Faraday rotation}

Faraday rotation affects the distribution of backscatter intensity across the different polarimetric channels. Its effects can be compensated if quad-pol data are available. If not available, it adds up to the error budget and reduces the performance. Therefore,  single-, dual- and compact-polarimetric configurations become potentially problematic. Fortunately, in many occasions the rotation angle is small enough to not distort severely the polarimetric information. 


\subsubsection{Input and Auxiliary Data}
\label{sec:aux_data_pol}

The input and auxiliary data necessary to generate the polarimetric products coincide with the \acs{nrb} product, as detailed in \autoref{tab:in_aux_pol_products}. 

\begin{longtblr}[
	label = tab:in_aux_pol_products,
	entry = none,
	caption = {\acs{cov}/\acs{gcov} input and auxiliary data},
	]{
		width = \linewidth,
		colspec = {Q[l,20]Q[l,60]Q[c,20]},
		hlines,
		hline{1-Z} = {-}{0.05em},
	}
	\textbf{Name}   & \textbf{Description}   & \textbf{Mandatory (M) / Optional (O)} \\
	\acs{slc} or \acs{cslc} & L1 Single Look Complex image or co-registered L1 Single Look Complex image & M \\
    \acs{slc} annotation & L1 Single Look Complex image annotations or reference geometry annotations & M \\
    \acs{dem} & External DEM (e.g., Copernicus DEM at different resolutions, with 30 m as the default worldwide). & M \\
    Orbits & \acs{slc} Standard, restituted or precise orbits & M \\
    \acs{etad} & \acs{etad}-like products for atmospheric delay correction & O \\
    \acs{tec} & External \acs{tec} data for delay correction & O \\
\end{longtblr}


\subsubsection{Main Variables}
\label{sec:variables_pol}

\autoref{tab:var_pol_products} lists the main variables associated with the \acs{cov}/\acs{gcov} products.

\begin{longtblr}[
	label = tab:var_pol_products,
	entry = none,
	caption = {\acs{cov}/\acs{gcov} main variables},
	]{
		width = \linewidth,
		colspec = {Q[l,15]Q[l,52]Q[l,20]Q[l,13]},
		hlines,
		hline{1-Z} = {-}{0.05em},
	}
	\textbf{Name} & \textbf{Description} & \textbf{Dimensions} & \textbf{Data Type} \\
	gamma0 & Radiometric Terrain Corrected Gamma-Naught image (backscatter in linear scale) of each channel & polarization, acquisition\_date, x, y & float32 \\

    correlation & Radiometric Terrain Corrected cross-correlation image (in linear scale) between pairs of channels & polarization pair, acquisition\_date, x, y & complex64 \\
        
    local\_inc
    
    \_angle & DEM-based local Incidence angle image [deg]. This variable is present only if the \acs{cov}/\acs{gcov} is generated from \acs{slc}s. If it is derived from \acs{cslc}s, this variable will be null, and the reference\_inc\_angle metadata will be present & x, y & float32 \\
    
    anf & Optional variable: area normalisation factor or local contributing area, used to obtain the gamma0 coefficient. This variable is present only if the \acs{cov}/\acs{gcov} is generated from \acs{slc}s. If it is derived from \acs{cslc}s, this variable will be null, and the reference\_anf metadata will be present & x, y & float32 \\
\end{longtblr}

The dimensions and coordinates associated with the \acs{gcov} main variables are described in \autoref{tab:var_gcov_coordinates}, whereas for \acs{cov} they are listed in \autoref{tab:var_cov_coordinates}. 

\begin{longtblr}[
	label = tab:var_gcov_coordinates,
	entry = none,
	caption = {\acs{gcov} main variables coordinates/dimensions},
	]{
		width = \linewidth,
		colspec = {Q[l,20]Q[l,62]Q[l,18]},
		hlines,
		hline{1-Z} = {-}{0.05em},
	}
	\textbf{Name} & \textbf{Description} & \textbf{Data Type} \\
	polarization & Polarimetric layers (entries of the matrix) & string \\
    acquisition\_date & Image acquisition dates, formatted as YYYY-MM-DD & string \\
    x & UTM/UPS x (East) coordinates [m] & float64 \\
    y & UTM/UPS y (North) coordinates [m] & float64 \\
\end{longtblr}

\begin{longtblr}[
	label = tab:var_cov_coordinates,
	entry = none,
	caption = {\acs{cov} main variables coordinates/dimensions},
	]{
		width = \linewidth,
		colspec = {Q[l,20]Q[l,62]Q[l,18]},
		hlines,
		hline{1-Z} = {-}{0.05em},
	}
	\textbf{Name} & \textbf{Description} & \textbf{Data Type} \\
	polarization & Polarimetric layers (entries of the matrix) & string \\
    acquisition\_date & Image acquisition dates, formatted as YYYY-MM-DD & string \\
    azimuth\_time & Time along satellite track direction & datetime64 \\
    range\_time & Time across satellite track direction & float64 \\
\end{longtblr}

\subsubsection{Main Metadata}
\label{sec:metadata_pol}

As it has been commented above, both \acs{cov} and \acs{gcov} products share most characteristics with the \acs{nrb} products described in~\autoref{s:nrb}, either because they are considered as an extension of the \acs{nrb} (as in the \acs{pol} product defined by \acs{ceos}) or because \acs{nrb} is a subset of \acs{gcov} (in the \acs{gcov} definition by NASA for NISAR). Consequently, all the product metadata described in~\autoref{sec:nrb_product_metadata} apply also to \acs{cov} and \acs{gcov}. 

However, for completeness and for making the search for information easier, all the product metadata are included in~\autoref{tab:meta_gcov_products}. Key comments about specific metadata which need special attention in the polarimetric \acs{ard} products are detailed. 

\autoref{tab:meta_gcov_products} lists the main metadata and per-pixel metadata associated with the \acs{gcov} product.

\begin{longtblr}[
	label = tab:meta_gcov_products,
	entry = none,
	caption = {\acs{gcov} main metadata and per-pixel metadata},
	]{
		width = \linewidth,
		colspec = {Q[l,17]Q[l,38]Q[l,20]Q[l,13]Q[c,12]},
		hlines,
		hline{1-Z} = {-}{0.05em},
	}
	\textbf{Name} & \textbf{Description} & \textbf{Dimensions} & \textbf{Data Type} & \textbf{Per-pixel Metadata}\\
	reference\_slc / reference\_cslc & List of links to input reference \acs{slc} or \acs{cslc} & - & list of string & \\
    reference\_dem & Link to reference DEM & - & string & \\
    etad\_file & ETAD file used for each image belonging to the \acs{cslc}/\acs{slc} stack. This is optional: it is included only if used; otherwise, it is null & - & string & \\ 
    tec & Link to external \acs{tec} data, if used & - & string & \\ 
    reference\_inc
    
    \_angle & Optional attribute: list of link to reference DEM-based local incidence angle maps. It is present only if the \acs{gcov} is generated from \acs{cslc}s. & - & string & \\
    reference\_anf & Optional attribute: list of links to reference area normalisation factor or local contributing area. It is present only if the \acs{gcov} is generated from \acs{cslc}s. & - & string & \\
    orig\_data\_id & Per pixel metadata reporting an index describing the bursts contributing to the output mosaicked pixel & x, y & int32 & \ding{51} \\
    mask & Mask containing, for each ground cell, the following information:
    
     - 0: no distortions
    
     - 1: layover
    
     - 2: shadow
    
     - 3: undefined & x, y & byte & \ding{51} \\
    nlooks & Number of equivalent looks for the gamma flattening & x, y & float32 & \ding{51} \\

    filtering & If speckle filtering has been applied, metadata should include: 
    
- Reference to algorithm

- Input filtering parameters

- Type

- Window size

- Other filter parameters & - & string & \\ 
    
    reference\_inc

    pol\_calib
    
    \_matrix 
    & Polarimetric distortion matrices with the channel imbalance and the cross-talk applied for the polarimetric calibration. (Optional) & - & complex64  &   \\

    mean\_Faraday
    
    \_rotation\_angle 
    
    &  Estimated mean Faraday rotation angle, obtained from the polarimetric data and/or from models, with link to the method or publication used to derive the estimate. (Optional) & - & float32  &  \\

   ionosphere
   
   \_mask 
   &  Flag indicating whether the image is “significantly impacted” by the ionosphere (0 – false, 1 – true). Significant impact would imply that the ionospheric impact on the backscatter exceeds the radiometric calibration requirement for the image. (Optional) & - & byte  &  \\
    
    geometric
    
    \_accuracy
    & Measured geometric accuracy, including eastern and northern geometric accuracy bias and standard deviation (only for projected products) & - & dictionary & \\ 
    east\_geom
    
    \_bias & Measured eastern geometric bias (only for projected products) & - & float32 & \\
    east\_geom\_std & Measured eastern geometric standard deviation (only for projected products) & - & float32 & \\
    north\_geom
    
    \_bias & Measured northern geometric bias (only for projected products) & - & float32 & \\
    north\_geom\_std & Measured northern geometric standard deviation (only for projected products) & - & float32 & \\
    geometric
    
    \_accuracy
    
    \_radial
    
    \_rmse	& Measured radial RMSE of geometric accuracy (only for projected products) & - & float32 & \\
\end{longtblr}

And \autoref{tab:var_pol_metadata_coordinates} describes the dimensions/coordinates associated with the \acs{cov} and \acs{gcov} main metadata.

\begin{longtblr}[
	label = tab:var_pol_metadata_coordinates,
	entry = none,
	caption = {\acs{cov} and \acs{gcov} main metadata coordinates/dimensions},
	]{
		width = \linewidth,
		colspec = {Q[l,20]Q[l,62]Q[l,18]},
		hlines,
		hline{1-Z} = {-}{0.05em},
	}
	\textbf{Name} & \textbf{Description} & \textbf{Data Type} \\
	x & UTM/UPS x (East) coordinates [m] & float64 \\
    y & UTM/UPS y (North) coordinates [m] & float64 \\
        azimuth\_time & Time along satellite track direction & datetime64 \\
    range\_time & Time across satellite track direction & float64 \\
\end{longtblr}

In the case of the \acs{cov} product, it contains the same metadata as \acs{gcov} because the products are ready for geocoding in the same way as the \acs{gcov} product. The only additional metadata is a look-up table with the correspondence between each output pixel in cartographical coordinates and the pixels in radar coordinates.


\subsubsection{Reference System Tiling and Slicing}
\label{sec:gcov_reference_system_tiling}

\autoref{tab:pol_ref_system_slicing} reports the reference system, the slicing, and grid for \acs{cov} and \acs{gcov} products.

\begin{longtblr}[
	label=tab:pol_ref_system_slicing,
	entry = none,
	caption = {Reference system and Slicing for \acs{gcov}},
	]{
		width=\linewidth,
		colspec = {Q[l,15]Q[l,35]Q[l,35]Q[l,15]},
		%vline{2-5} = {-}{},
		%vline{5-5} = {-}{},
		hlines,
		hline{1-Z} = {-}{0.05em},
        rowhead=2
	}
		\textbf{Product} & \textbf{Reference System} & \textbf{Slicing} & \textbf{Spacing}\\
	\textbf{\acs{cov}} & Radar: azimuth/slant-range & Burst & Same of the input data \\
	\textbf{\acs{gcov}} & Cartographic: UTM/UPS & 100 km x 100 km & 10 m\\   
\end{longtblr}

In \autoref{tab:pol_ref_system_slicing}, the slicing for the \acs{cov} product is assigned as a burst as it is derived from \acs{slc} or \acs{cslc} products. However, the correspondence data necessary to convert it to \acs{gcov}, e.g. a look-up table with the association of radar coordinates corresponding to each cartographical coordinate, should be included as a metadata layer. After this conversion, the resulting product is really an \acs{ard} with total integrability and interoperability, like the \acs{gcov} product.

\subsection{Product Requirements}
\label{sec:impacts_pol}

The main requirements for a correct generation of \acs{cov} and \acs{gcov} products coincide with the ones stated for \acs{nrb} in \autoref{sec:impacts_nrb}, and the processing steps required to generate them coincide also with all the steps carried out to obtain \acs{nrb} products. 

All the processing steps must be applied in the same way (same algorithm and same parameters) to each individual entry of the polarimetric covariance matrix. It must be noted that the off-diagonal elements of the polarimetric matrix are complex numbers (not real numbers as the diagonal elements),  but they must be processed in the same way as the diagonal elements.



\subsubsection{Accuracy}

\begin{longtblr}[
	label = tab:pol_accuracy_cband,
	entry = none,
	caption = {Accuracy Requirements for \acs{cov} and \acs{gcov} products (C-Band)},
	]{
		width = \linewidth,
		colspec = {Q[l,30] Q[l,60] Q[c,20]},
		hlines,
		hline{1-Z} = {-}{0.05em},
		rowhead = 1,
	}
	
	\textbf{Metric} & \textbf{Description} & \textbf{Requirement} \\
	
	\SetCell[r=2]{l} Range Accuracy (\%) 
	& Model-Driven (only orbit and \acs{dem}) & $\leq 5\%$\\
	& Model-Driven + atmospheric correction & $\leq 2\%$ \\
	
	\SetCell[r=2]{l} Azimuth Accuracy (\%) 
	& Model-Driven (only orbit and \acs{dem}) & $\leq 2\%$\\
	& Data-Driven (\acs{esd}) & $\leq 0.5\%$ \\
	
	\acs{rcs} difference (dB) & Stability of the \acs{rcs} comparing different acquisition geometries over the same area considering a resolution of 30m & $\leq 5$~dB \\
	
\end{longtblr}

\begin{longtblr}[
	label = tab:pol_accuracy_lband,
	entry = none,
	caption = {Accuracy Requirements for \acs{cov} and \acs{gcov} products (L-Band)},
	]{
		width = \linewidth,
		colspec = {Q[l,30] Q[l,60] Q[c,20]},
		hlines,
		hline{1-Z} = {-}{0.05em},
		rowhead = 1,
	}
	
	\textbf{Metric} & \textbf{Description} & \textbf{Requirement} \\
	
	\SetCell[r=3]{l} Range Accuracy (\%) 
	& Model-Driven (only orbit and \acs{dem}) & $\leq 60\%$\\
	& Model-Driven + atmospheric correction & $\leq 25\%$ \\
	& Model-Driven + atmospheric correction + Data-Driven (\acs{icc}) & $\leq 10\%$ \\
	
	\SetCell[r=3]{l} Azimuth Accuracy (\%) 
	& Model-Driven (only orbit and \acs{dem}) & $\leq 150\%$ \\
    & Model-Driven + atmospheric correction & $\leq 50\%$ \\
	& Data-Driven (\acs{icc}) & $\leq 10\%$ \\
	
	\acs{rcs} difference (dB) & Stability of the \acs{rcs} comparing different acquisition geometries over the same area considering a resolution of 30m & $\leq 5$~dB \\
	
\end{longtblr}


\subsubsection{Timeliness}

The \acs{cov} and \acs{gcov} products can be generated as soon as L1 data is available for each \acs{slc} image independently or starting from \acs{cslc} products. The generation of covariance matrix for the provision of the polarimetric products does not impact significantly the timeline. 

Besides choosing between \acs{slc} or \acs{cslc} products, the main factor impacting the timeline is the availability of external data or reference images, which are required to refine geolocation and mitigate atmospheric effects, particularly ionospheric effects. 


\autoref{fig:pol_single} shows the flowchart of generating the \acs{cov} and \acs{gcov} products using a single \acs{slc} image. The associated timeline was provided in~\autoref{fig:nrb_single}, indicating the time required to collect the input data and complete the main processing steps. 

\begin{figure}[H]
	\centering
	\resizebox{\textwidth}{!}{
		\begin{tikzpicture}[node distance=1cm and 2cm]
			
			% Define the style for the nodes
			% \tikzstyle{node} = [circle, draw, fill=blue!20, text centered, minimum size=1.5cm]
			
			% Draw nodes
			\node[rectin] (init) {\acs{slc}};
			\node[block, right=of init] (g2r) {Geo-to-Radar};
            \node[block, right=of g2r] (rtc) {\acs{rtc}};
			\node[block, right=of rtc] (covm) {Covariance Matrix};
			\node[block, right=of covm] (geocoding) {Geocoding};
			\node[rectaux, below=of init] (dem) {\acs{dem}};
			\node[rectaux, below=of dem] (orbit) {Orbital data};
			\node[rectauxd, below=of orbit] (models) {External models};
			\node[rectind, below=of models] (ref_stack) {Reference SLC stack};
			\node[rectout, right=of geocoding] (gcov) {\acs{gcov}};
			\node[rectout, below=of covm] (cov) {\acs{cov}};
			
			% Draw edges with arrows
			\draw[->, thick] (init) -- (g2r);
			\draw[->, thick] (g2r) -- (rtc);
			\draw[->, thick] (rtc) --(covm);
			\draw[->, thick] (covm) -- (geocoding);
			\draw[->, thick] (covm) -- (cov);
			\draw[->, thick] (dem) -| (g2r);
			\draw[->, thick] (orbit) -| (g2r);
			\draw[->, thick] (models) -| (g2r);
			\draw[->, thick] (ref_stack) -| (g2r);
			\draw[->, thick] (geocoding) -- (gcov);
		\end{tikzpicture}
	}
	\caption{Sample workflow for generation of \acs{cov} and \acs{gcov} considering each \acs{slc} image separately. The step denoted as "Covariance Matrix" should include a speckle filtering, but is declared as optional. In the \acs{gcov} product defined by NISAR the speckle filtering and the geocoding are carried out simultaneously with an area-adaptive algorithm. The associated timeline can be consulted in the corresponding figure for \acs{nrb} (\autoref{fig:nrb_single}).}
	\label{fig:pol_single}
\end{figure}


In case the \acs{cov} or \acs{gcov} products are generated from the \acs{cslc}, all the calibrations have already been performed during the \acs{cslc} processing, and the reference geometry is computed only once for the reference orbit. Therefore, the schema simplifies as reported in \autoref{fig:pol_from_cslc}, but the timeline does not change depending on the information needed to perform Geo-to-Radar conversion or to generate the \acs{cslc} product.
\begin{figure}[H]
	\centering
	\resizebox{\textwidth}{!}{
		\begin{tikzpicture}[node distance=0.5cm and 2cm]
			
			% Define the style for the nodes
			% \tikzstyle{node} = [circle, draw, fill=blue!20, text centered, minimum size=1.5cm]
			
			% Draw nodes
			\node[rectin] (init) {\acs{slc}};
			\node[block, right=of init] (coreg) {Co-registration};
			\node[rectout, right=of coreg] (cslc) {\acs{cslc}};
            \node[block, right=of cslc] (rtc) {\acs{rtc}};
			\node[block, right=of rtc] (covm) {Covariance Matrix};
			\node[block, right=of covm] (geocoding) {Geocoding};
			\node[rectout, right=of geocoding] (gcov) {\acs{gcov}};
			\node[rectout, below=of covm] (cov) {\acs{cov}};
			
			\node[rectaux, below=of init] (dem) {\acs{dem}};
			\node[rectaux, below=of dem] (orbit) {Orbital data};
			\node[rectauxd, below=of orbit] (models) {External models};
			\node[rectind, below=of models] (ref_stack) {Reference SLC stack};
			
			% Draw edges with arrows
			\draw[->, thick] (init) -- (coreg);
			\draw[->, thick] (coreg) -- (cslc);
			\draw[->, thick] (cslc) -- (rtc);
			\draw[->, thick] (rtc) -- (covm);
			\draw[->, thick] (covm) -- (cov);
			\draw[->, thick] (covm) -- (geocoding);
			\draw[->, thick] (geocoding) -- (gcov);
			
			% Draw edges with arrows
			\draw[->, thick] (dem) -| (coreg);
			\draw[->, thick] (orbit) -| (coreg);
			\draw[->, thick] (models) -| (coreg);
			\draw[->, thick] (ref_stack) -| (coreg);
			
		\end{tikzpicture}
	}
	\caption{Sample workflow for generation of \acs{cov} and \acs{gcov} starting from \acs{cslc} products. The step denoted as "Covariance Matrix" should include a speckle filtering, but is declared as optional. In the \acs{gcov} product defined by NISAR the speckle filtering and the geocoding are carried out simultaneously with an area-adaptive algorithm. The associated timeline can be consulted in the corresponding figure for \acs{nrb} (\autoref{fig:nrb_from_cslc}).}
	\label{fig:pol_from_cslc}
\end{figure}

\subsection{Usability, Integrability and Interoperability}
\label{sec:usability_pol}

The characteristics of the \acs{gcov} products in terms of usability, integrability and interoperability are coincident, and extended, with respect to the \acs{nrb} products, which were detailed in \autoref{tab:nrb_usability_integrability_interoperability}. This is so because the first obvious way to use polarimetric products is by considering only the diagonal values of the covariance matrix. This strategy is equivalent to working with a \acs{nrb} product with several layers of intensity data (one layer per polarimetric channel). As such, all the comments about the usage of \acs{nrb} products (see \autoref{sec:nrb_usability_integrability}) apply here. 

However, it is when the matrix form is considered as a joint feature or entity when polarimetry really excels in many application domains. Therefore, the specific usability, integrability and interoperability of these matrix-form products is described next in \autoref{tab:gcov_usability_integrability_interoperability}. 

\begin{longtblr}[
  caption = {Usability, Integrability, and Interoperability of \acs{gcov} product},
  label = {tab:gcov_usability_integrability_interoperability},
]{
  colspec = {Q[2cm] Q[12cm]},
  rowhead = 1,
  width = \linewidth
}
\hline
\textbf{Aspect} & \textbf{Description} \\
\hline
Usability &
\begin{itemize}
  \item Radiometric calibration to $\gamma^0$ convention enables direct comparison across sensors, acquisition times, and geometries.
  \item Orthorectified and georeferenced product makes data ready to use in thematic mapping applications (classification and change detection). 
  \item Polarimetric features: Besides the direct use of the entries of the covariance matrix, many polarimetric features can be derived from them. The most clear example is the outputs of polarimetric decompositions~\cite{Lee_Pottier2009}. Starting from \acs{gcov} products, these features can be readily computed.
  \item Machine Learning: \acs{gcov} products, as well as decomposition outputs, can be directly ingested into machine learning models for classification, regression, or segmentation tasks without the need for extra processing. The input features to these models are usually expressed as layers or bands of data, which in this case would be formed by the entries of the covariance matrix and the outputs of decompositions. In the case the machine learning algorithm is prepared only for intensity data, there exist simple expressions which allow bijectively converting the complex matrix into a set of intensities, and vice versa~\cite{Hoekman2003, Mestre2024}.
\end{itemize}
\\

Integrability &
\begin{itemize}
  \item GIS Compatibility: Easily used within \acs{gis} platforms, due to geographic referencing and data format.
  \item Multi-Sensor Fusion: Integrates well with optical or other SAR projected products for land cover, surface change, or hazard mapping.
  \item Environmental Monitoring: Supports use in models or indicators for forest cover, crop health, or hydrology. 
  \item Users can overlay the matrix elements contained in these products with vector layers, other remote sensing data (optical, infrared), or ground truth data for comprehensive spatial analysis.
\end{itemize}
\\

Interoperability &
\begin{itemize}
  \item Cross-Mission Compatibility: Normalized $\gamma^0$ used in \acs{gcov} allows for comparison between satellites.
  \item Interdisciplinary Applications: Usable in hydrology, agriculture, forestry, urban studies, etc., enhancing cross-domain use.
  \item Compatible with common \acs{eo} frameworks and pipelines using georeferenced data.
  \item High interoperability due to standard geographic referencing and pixel alignment across scenes.
  \item Adheres to ARD principles, enhancing integration with downstream services and platforms
  \item Facilitates cross-mission usage through harmonized ground projection.
\end{itemize}
\\
\hline
\end{longtblr}



